%%
%% This is a placeholder manuscript shell for WWW'27 (Research Track — Security and Privacy).
%% Refresh this manuscript before final submission by syncing with docs/paper/manuscript.tex
%% from the research repository (demo-21-mcp-ecosystem-measurement-eurosp).
%%
\documentclass[sigconf,anonymous,review]{acmart}

\acmConference[WWW '27]{The ACM Web Conference 2027}{May 10--14, 2027}{Dublin, Ireland}
\acmYear{2027}
\acmDOI{}
\acmISBN{}

\title{Connecting Agents to Web Services: HTTP Endpoints, Secret Headers and Version Changes in the Model Context Protocol Registry}

\author{Anonymous Author(s)}
\affiliation{%
  \institution{Anonymous Institution}
  \country{Anonymous Country}
}
\email{anonymous@example.org}

\begin{abstract}
The official Model Context Protocol (MCP) registry walk of 30 September 2026 contains 126,847 version records across 37,815 servers, of which 61.89\% declare remote HTTP endpoints and 56.96\% operate purely as remote services without installable packages. This census measures the emerging protocol connecting autonomous agents to web services. Remote endpoints resolve to 15,738 distinct hosts across 6,432 registrable domains, with 79.5\% unranked in the Tranco Top-1M. Furthermore, 3,679 remotes declare confidential headers marked \texttt{isSecret}, binding ambient credentials to outbound requests. Across 89,032 publication-ordered version transitions, releases arrive rapidly (median 0.8 days, 52.4\% within 24 hours). Transitions modify remote URLs in 1,353 cases: 52.4\% alter destination hostnames, 36.1\% cross administrative eTLD+1 boundaries, and 69.8\% of cross-domain migrations target multi-tenant serverless or PaaS platforms. Evaluating a deterministic pattern screen over 40,672 tool descriptions yields a 1.25\% flag rate; on MCPTox attacks, it achieves only 6.2\% recall while flagging benign documentation in 89.2\% of real-world hits. Rather than claiming unverified attack prevalence, the paper quantifies longitudinal structural shifts and proposes three architectural defences---origin-bound credential scoping, manifest attestations, and domain lifecycle verification---to secure agentic web interfaces.
\end{abstract}

\begin{CCSXML}
<ccs2012>
   <concept>
       <concept_id>10002951.10003260.10003282</concept_id>
       <concept_desc>Information systems~Web applications</concept_desc>
       <concept_significance>500</concept_significance>
   </concept>
   <concept>
       <concept_id>10002978.10003006</concept_id>
       <concept_desc>Security and privacy~Systems security</concept_desc>
       <concept_significance>300</concept_significance>
   </concept>
   <concept>
       <concept_id>10002951.10003260.10003277</concept_id>
       <concept_desc>Information systems~Web mining</concept_desc>
       <concept_significance>300</concept_significance>
   </concept>
</ccs2012>
\end{CCSXML}

\ccsdesc[500]{Information systems~Web applications}
\ccsdesc[300]{Security and privacy~Systems security}
\ccsdesc[300]{Information systems~Web mining}

\keywords{Model Context Protocol, agentic systems, Web measurement, tool poisoning, HTTP endpoints, provenance attestations}

\begin{document}

\maketitle

\section{Introduction}
As autonomous language model agents transition from sandboxed code interpreters to distributed web actors, they increasingly interact with external web APIs, cloud microservices, and databases through standardised tool-calling protocols. The Model Context Protocol (MCP) has rapidly emerged as a dominant open standard for this interface: an MCP server is addressed by a remote HTTP endpoint or packaged as an installable binary, discovered through a public registry that operates as a Web API, and described to the model via natural-language text.

This paper presents a longitudinal measurement study of the official MCP registry walk, focusing on Web infrastructure, transport distribution, secret-header configurations, and URL mutations across releases.

\section{Measurement Methodology}
The empirical analysis is based on secondary data collected from the official registry walk of 30 September 2026, comprising 126,847 version records across 37,815 distinct servers. All analyses run deterministically from committed pipeline scripts.

\section{Findings}
Over 61.89\% of registered servers declare remote HTTP endpoints. Longitudinal tracking reveals 1,353 transitions modifying remote URLs, with 52.4\% altering destination hostnames and 36.1\% crossing administrative eTLD+1 boundaries.

\section{Defences and Architectural Recommendations}
We propose origin-bound credential scoping, out-of-band manifest attestations, and registry domain lifecycle verification to harden the agentic Web interface.

\bibliographystyle{ACM-Reference-Format}
\bibliography{references}

\end{document}
